\hypertarget{appendix-0234ab7f86b1}{}\label{referencestart:appendix-0234ab7f86b1}
\section{Subdivision surface masks}\label{appendix-0234ab7f86b1-subdivision-surface-masks}

\hypertarget{appendix-0234ab7f86b1-subdivision-surface-masks}{}

\textbf{Curriculum depth:} Developed · \textbf{First seen:} University courses · \textbf{Latest development:} University courses

\begin{allmightroledefinition}

\subsection{Definition and mechanism}\label{appendix-0234ab7f86b1-definition-and-mechanism}

\hypertarget{appendix-0234ab7f86b1-definition-and-mechanism}{}

Subdivision refines a control mesh and updates old and newly inserted vertices with local masks, producing a sequence that approaches a smooth limit surface. Loop subdivision targets triangle meshes: for an interior edge with endpoints \(a,b\) and opposite vertices \(c,d\), its new point is \(3(a+b)/8+(c+d)/8\). An old vertex of degree \(n\) becomes \((1-n\beta)S+\beta\sum_iP_i\), with the lecture's simple weight \(\beta=3/16\) for \(n=3\) and \(3/(8n)\) otherwise. Catmull--Clark targets polygon meshes: face points are face-vertex averages, interior edge points average endpoints and two adjacent face points, and old vertices use face-average, edge-midpoint-average, and original-position terms. Boundary rules differ from interior masks and must be stated. Extraordinary vertices have nonregular valence and need deliberate treatment.

\end{allmightroledefinition}

\subsection{Prerequisites}\label{appendix-0234ab7f86b1-prerequisites}

\hypertarget{appendix-0234ab7f86b1-prerequisites}{}

\begin{itemize}
\tightlist
\item
  \hyperref[appendix-1aea26b2307f]{Half-edge mesh connectivity}
\item
  \href{../07-Vision-XR-and-Playable-Worlds-P01/main.pdf\#appendix-0026ec42c9d5}{Convex combinations and convex hulls (Vision, XR, and Playable Worlds, Part 1, Reference notes: Convex combinations and convex hulls)}
\end{itemize}

\subsection{Used by}\label{appendix-0234ab7f86b1-used-by}

\hypertarget{appendix-0234ab7f86b1-used-by}{}

\begin{itemize}
\item
  {Worlds, Images and Inference}
\item
  \href{../03-Geometry-Cameras-and-Image-Formation-P01/main.pdf\#appendix-7107009c505b}{Local rigidity and mesh orientation fields (Geometry, Cameras, and Image Formation, Part 1, Reference notes: Local rigidity and mesh orientation fields)} \#\# Full teaching explanation
\end{itemize}

\href{../03-Geometry-Cameras-and-Image-Formation-P01/main.pdf\#03-visual-computing-i-ch002-meshes-connectivity-and-geometry-processing}{Meshes, connectivity, and geometry processing (Geometry, Cameras, and Image Formation, Part 1, Ch. 9)} · \href{../07-Vision-XR-and-Playable-Worlds-P01/main.pdf\#04-visual-computing-ii-ch032-surface-descriptions-and-gaussian-images}{Surface descriptions and Gaussian images (Vision, XR, and Playable Worlds, Part 1, Ch. 3)}
